\documentclass[11pt]{article}

% ============================================================
% Packages
% ============================================================
\usepackage[utf8]{inputenc}
\usepackage[margin=1in]{geometry}
\usepackage{amsmath,amssymb,amsthm}
\usepackage{graphicx}
\usepackage{booktabs}
\usepackage{xcolor}
\usepackage{hyperref}
\usepackage{tikz}
\usepackage{enumitem}
\usepackage{pgfplots}

\pgfplotsset{compat=1.18}

% ============================================================
% Document Formatting
% ============================================================
\setlength{\parskip}{0.6em}
\setlength{\parindent}{0pt}

% ============================================================
% List Formatting
% ============================================================
\setlist[enumerate]{
    leftmargin=2em,
    itemsep=0.25em
}

% ============================================================
% Solution Environment
% ============================================================
\newenvironment{solution}
    {\par\noindent\textbf{Solution.}\par}
    {\par}

% ============================================================
% Link Styling
% ============================================================
\hypersetup{
    colorlinks=true,
    linkcolor=blue,
    urlcolor=cyan
}

% ============================================================
% Metadata
% ============================================================
\title{Ordinary Differential Equations}
\author{Chris Ritter}
\date{August 20, 2026}

% ============================================================
% Document
% ============================================================

\begin{document}

\maketitle
\tableofcontents
\newpage

% ============================================================
\section{Integral Curves}
% ============================================================

An \textbf{ordinary differential equation (ODE)} describes a
relationship between an unknown function and one or more of its
derivatives.

For example,

\[
\frac{dy}{dt}=ay-b,
\]

where $a$ and $b$ are constants.

A \textbf{solution} to a differential equation is a function that
satisfies the equation. The graph of a solution is called an
\textbf{integral curve}.

% ------------------------------------------------------------
\subsection{A Population Model}
% ------------------------------------------------------------

Consider the population model

\[
\frac{dp}{dt}=rp-k,
\]

where $p(t)$ represents the population at time $t$, $r$ is a
proportional growth rate, and $k$ is a constant rate of population
loss.

Suppose

\[
r=\frac{1}{2},
\qquad
k=450.
\]

Then

\[
\frac{dp}{dt}
=
\frac{1}{2}p-450.
\]

Factoring,

\[
\frac{dp}{dt}
=
\frac{1}{2}(p-900).
\]

% ------------------------------------------------------------
\subsection{Separation of Variables}
% ------------------------------------------------------------

To solve the equation, separate the variables:

\[
\frac{dp}{p-900}
=
\frac{1}{2}\,dt.
\]

Integrating both sides,

\[
\int \frac{dp}{p-900}
=
\int \frac{1}{2}\,dt.
\]

Therefore,

\[
\ln|p-900|
=
\frac{t}{2}+C.
\]

Exponentiating,

\[
|p-900|
=
e^{t/2+C}.
\]

Since the constant can absorb the sign, we write

\[
p-900
=
\alpha e^{t/2},
\]

where $\alpha$ is an arbitrary constant.

Thus,

\[
\boxed{
p(t)=900+\alpha e^{t/2}
}.
\]

% ------------------------------------------------------------
\subsection{Applying an Initial Condition}
% ------------------------------------------------------------

Suppose the initial population is

\[
p(0)=p_0.
\]

Substituting $t=0$,

\[
p_0
=
900+\alpha e^0.
\]

Since $e^0=1$,

\[
p_0=900+\alpha.
\]

Therefore,

\[
\alpha=p_0-900.
\]

The solution satisfying the initial condition is

\[
\boxed{
p(t)
=
900+(p_0-900)e^{t/2}
}.
\]

% ------------------------------------------------------------
\subsection{Equilibrium Population}
% ------------------------------------------------------------

An equilibrium occurs when the population does not change. Thus,

\[
\frac{dp}{dt}=0.
\]

For the population model,

\[
\frac{1}{2}p-450=0.
\]

Solving,

\[
p=900.
\]

Therefore,

\[
\boxed{p=900}
\]

is the equilibrium population.

Notice that the general solution

\[
p(t)=900+\alpha e^{t/2}
\]

contains the equilibrium solution when $\alpha=0$.

% ============================================================
\section{The General Equation \texorpdfstring{$dy/dt=ay-b$}{dy/dt=ay-b}}
% ============================================================

Consider the differential equation

\[
\frac{dy}{dt}=ay-b,
\]

where $a$ and $b$ are constants.

Factor out $a$:

\[
\frac{dy}{dt}
=
a\left(y-\frac{b}{a}\right).
\]

Assuming $a\neq0$, separate the variables:

\[
\frac{dy}{y-b/a}
=
a\,dt.
\]

Integrating both sides,

\[
\int \frac{dy}{y-b/a}
=
\int a\,dt.
\]

Therefore,

\[
\ln\left|y-\frac{b}{a}\right|
=
at+C.
\]

Exponentiating,

\[
y-\frac{b}{a}
=
\alpha e^{at}.
\]

Thus,

\[
\boxed{
y(t)
=
\frac{b}{a}
+
\alpha e^{at}
}.
\]

% ------------------------------------------------------------
\subsection{Applying an Initial Condition}
% ------------------------------------------------------------

Suppose

\[
y(0)=y_0.
\]

Then

\[
y_0
=
\frac{b}{a}
+
\alpha e^0.
\]

Since $e^0=1$,

\[
y_0
=
\frac{b}{a}
+
\alpha.
\]

Therefore,

\[
\alpha
=
y_0-\frac{b}{a}.
\]

Hence,

\[
\boxed{
y(t)
=
\frac{b}{a}
+
\left(
y_0-\frac{b}{a}
\right)e^{at}
}.
\]

% ============================================================
\section{Differential Equations and Initial-Value Problems}
% ============================================================

It is important to distinguish between a differential equation and an
initial-value problem.

A \textbf{differential equation} specifies a relationship involving an
unknown function and its derivatives.

For example,

\[
\frac{dy}{dt}=ay-b.
\]

The general solution is

\[
y(t)
=
\frac{b}{a}
+
\alpha e^{at}.
\]

This represents a family of solutions, with each value of $\alpha$
giving a different solution.

An \textbf{initial condition} selects one particular solution from this
family. For example,

\[
y(0)=y_0.
\]

Therefore,

\[
\boxed{
\text{Differential Equation}
+
\text{Initial Condition}
=
\text{Initial-Value Problem (IVP)}
}.
\]

% ============================================================
\section{Population Growth and Decay}
% ============================================================

A common population model is

\[
\frac{dp}{dt}=rp-k.
\]

Here,

\begin{itemize}
    \item $p(t)$ is the population at time $t$,
    \item $r$ is the proportional growth rate,
    \item $k$ is a constant rate of population loss.
\end{itemize}

The equilibrium population occurs when

\[
\frac{dp}{dt}=0.
\]

Therefore,

\[
rp-k=0,
\]

which gives

\[
p=\frac{k}{r}.
\]

Thus, the equilibrium population is

\[
\boxed{
p_{\mathrm{eq}}=\frac{k}{r}
}.
\]

The general solution is

\[
\boxed{
p(t)
=
\frac{k}{r}
+
\left(
p_0-\frac{k}{r}
\right)e^{rt}
}.
\]

% ------------------------------------------------------------
\subsection{Interpretation of the Solution}
% ------------------------------------------------------------

The solution

\[
p(t)
=
\frac{k}{r}
+
\left(
p_0-\frac{k}{r}
\right)e^{rt}
\]

can be understood by comparing the initial population $p_0$ with the
equilibrium population $k/r$.

If

\[
p_0=\frac{k}{r},
\]

then

\[
p(t)=\frac{k}{r}
\]

for all $t$.

If

\[
p_0>\frac{k}{r},
\]

then the initial population lies above equilibrium.

If

\[
p_0<\frac{k}{r},
\]

then the initial population lies below equilibrium.

% ============================================================
\section{Falling Objects}
% ============================================================

A falling object provides another important application of first-order
differential equations.

Suppose an object has mass $m$ and is subject to gravity and air
resistance. Assume that the drag force is proportional to velocity.

The model is

\[
m\frac{dv}{dt}=mg-\gamma v,
\]

where

\begin{itemize}
    \item $m$ is the mass of the object,
    \item $g$ is the acceleration due to gravity,
    \item $v(t)$ is the velocity,
    \item $\gamma>0$ is the drag coefficient.
\end{itemize}

Dividing by $m$ gives

\[
\frac{dv}{dt}
=
g-\frac{\gamma}{m}v.
\]

This has the same basic form as the previous equations.

% ------------------------------------------------------------
\subsection{Solving the Falling-Object Equation}
% ------------------------------------------------------------

Rewrite the equation as

\[
\frac{dv}{dt}
=
-\frac{\gamma}{m}
\left(
v-\frac{mg}{\gamma}
\right).
\]

Separate the variables:

\[
\frac{dv}
{v-mg/\gamma}
=
-\frac{\gamma}{m}\,dt.
\]

Integrating,

\[
\ln\left|
v-\frac{mg}{\gamma}
\right|
=
-\frac{\gamma}{m}t+C.
\]

Exponentiating,

\[
v-\frac{mg}{\gamma}
=
\alpha e^{-\gamma t/m}.
\]

Therefore,

\[
v(t)
=
\frac{mg}{\gamma}
+
\alpha e^{-\gamma t/m}.
\]

If the initial velocity is

\[
v(0)=v_0,
\]

then

\[
v_0
=
\frac{mg}{\gamma}
+
\alpha.
\]

Thus,

\[
\alpha
=
v_0-\frac{mg}{\gamma}.
\]

The solution is therefore

\[
\boxed{
v(t)
=
\frac{mg}{\gamma}
+
\left(
v_0-\frac{mg}{\gamma}
\right)
e^{-\gamma t/m}
}.
\]

% ------------------------------------------------------------
\subsection{Terminal Velocity}
% ------------------------------------------------------------

The quantity

\[
\boxed{
v_T=\frac{mg}{\gamma}
}
\]

is called the \textbf{terminal velocity}.

As $t\to\infty$,

\[
e^{-\gamma t/m}\longrightarrow0.
\]

Therefore,

\[
\lim_{t\to\infty}v(t)
=
\frac{mg}{\gamma}.
\]

Hence,

\[
\boxed{
v(t)\longrightarrow\frac{mg}{\gamma}
\qquad\text{as }t\to\infty.
}
\]

Thus, the velocity approaches a limiting value rather than increasing
without bound.

% ------------------------------------------------------------
\subsection{Example: Falling Object}
% ------------------------------------------------------------

Consider a falling object with

\[
g=9.8,
\qquad
m=10,
\qquad
\gamma=2.
\]

Assume the object is initially at rest, so

\[
v(0)=0.
\]

The velocity function is

\[
v(t)
=
\frac{mg}{\gamma}
+
\left(
v_0-\frac{mg}{\gamma}
\right)e^{-\gamma t/m}.
\]

Substituting the given values,

\[
v(t)
=
\frac{(10)(9.8)}{2}
+
\left(
0-\frac{(10)(9.8)}{2}
\right)e^{-(2/10)t}.
\]

Therefore,

\[
\boxed{
v(t)=49-49e^{-0.2t}
}.
\]

% ------------------------------------------------------------
\subsubsection{Finding the Time to Reach the Ground}
% ------------------------------------------------------------

Suppose the object is dropped from a height of $300$ meters.

Let $s(t)$ denote the distance the object has fallen. Since velocity is
the derivative of position,

\[
\frac{ds}{dt}=v(t).
\]

Thus,

\[
\frac{ds}{dt}
=
49-49e^{-0.2t},
\qquad
s(0)=0.
\]

Integrating,

\[
\int ds
=
\int
\left(
49-49e^{-0.2t}
\right)dt.
\]

Therefore,

\[
s(t)
=
49t+245e^{-0.2t}+C.
\]

Using $s(0)=0$,

\[
0
=
49(0)+245e^0+C.
\]

Hence,

\[
C=-245.
\]

Therefore,

\[
\boxed{
s(t)=49t+245e^{-0.2t}-245
}.
\]

The object reaches the ground when

\[
s(t^*)=300.
\]

Thus,

\[
49t^*
+
245e^{-0.2t^*}
-
245
=
300.
\]

This equation must be solved numerically. The solution is approximately

\[
\boxed{
t^*\approx10.51\text{ seconds}.
}
\]

% ============================================================
\section{Classification of Differential Equations}
% ============================================================

Differential equations can be classified according to several
different properties.

% ------------------------------------------------------------
\subsection{A. Number of Independent Variables}
% ------------------------------------------------------------

A differential equation is called an \textbf{ordinary differential
equation (ODE)} if there is only one independent variable.

For example,

\[
\frac{dv}{dt}
=
g-\frac{\gamma}{m}v
\]

is an ODE because the only independent variable is $t$.

A differential equation with more than one independent variable is
called a \textbf{partial differential equation (PDE)}.

For example, the heat equation can be written as

\[
\boxed{
\frac{\partial u}{\partial t}
=
\alpha^2
\frac{\partial^2u}{\partial x^2}
}.
\]

Here, $u=u(x,t)$ depends on both $x$ and $t$.

% ------------------------------------------------------------
\subsection{B. Number of Unknown Functions}
% ------------------------------------------------------------

A differential equation may involve a single unknown function or
multiple unknown functions.

A collection of differential equations involving multiple unknown
functions is called a \textbf{system of differential equations}.

For example,

\[
\frac{dx}{dt}
=
\alpha x-\beta xy,
\]

\[
\frac{dy}{dt}
=
\gamma xy-\theta y.
\]

This is a system involving the two unknown functions $x(t)$ and $y(t)$.

% ------------------------------------------------------------
\subsection{C. Order of a Differential Equation}
% ------------------------------------------------------------

The \textbf{order} of a differential equation is the order of the
highest derivative that appears in the equation.

In general, an $n$th-order ODE can be written as

\[
F\left(
t,y,y',y'',\ldots,y^{(n)}
\right)=0.
\]

An equation can also be written explicitly in terms of its highest
derivative:

\[
y^{(n)}
=
f\left(
t,y,y',\ldots,y^{(n-1)}
\right).
\]

For example,

\[
y^{(3)}
+
2e^t y''
+
y'
-
t^4
=
0
\]

is a third-order differential equation because the highest derivative
is $y^{(3)}$.

Another example is

\[
(y')^2+ty'+4y=0.
\]

This is a first-order differential equation because the highest
derivative appearing is $y'$.

Solving this equation for $y'$ gives

\[
(y')^2+ty'+4y=0.
\]

Using the quadratic formula,

\[
\boxed{
y'
=
\frac{-t\pm\sqrt{t^2-16y}}{2}
}.
\]

% ------------------------------------------------------------
\subsection{D. Linear vs. Nonlinear}
% ------------------------------------------------------------

An $n$th-order differential equation

\[
F\left(
t,y,y',\ldots,y^{(n)}
\right)=0
\]

is \textbf{linear} if it is linear in

\[
y,\quad y',\quad y'',\quad\ldots,\quad y^{(n)}.
\]

The most general form of an $n$th-order linear differential equation is

\[
\boxed{
a_0(t)y^{(n)}
+
a_1(t)y^{(n-1)}
+
\cdots
+
a_{n-1}(t)y'
+
a_n(t)y
=
g(t)
}.
\]

The functions

\[
a_0(t),a_1(t),\ldots,a_n(t)
\]

are coefficients that depend only on the independent variable $t$.

% ------------------------------------------------------------
\subsubsection{Example: Linear Equation}
% ------------------------------------------------------------

Consider

\[
\frac{t}{1+t^2}y''
+
3ty
=
\sin(t).
\]

This equation is

\begin{itemize}
    \item ordinary, because there is one independent variable $t$;
    \item second-order, because the highest derivative is $y''$;
    \item linear, because $y''$ and $y$ appear only to the first power
    and are not multiplied together.
\end{itemize}

Therefore, this is an

\[
\boxed{\text{ordinary, linear, second-order ODE}.}
\]

% ------------------------------------------------------------
\subsubsection{Examples of Nonlinear Equations}
% ------------------------------------------------------------

Consider

\[
yy'=0.
\]

This equation is \textbf{nonlinear} because the unknown function $y$
and its derivative $y'$ are multiplied together.

Similarly,

\[
y^{(3)}
+
2e^t y''
+
yy'
=
t^4
\]

is nonlinear because of the product

\[
yy'.
\]

Therefore,

\[
\boxed{
y^{(3)}+2e^t y''+yy'=t^4
}
\]

is a third-order nonlinear ODE.

\end{document}
